\section{Diagrams and their embedding tests}
\label{sec:maps-elementarity-04}
\label{sec:diagrams}

\begin{definition}[Constants for parameters]
\label{def:parameter-expansion}
If \(A\subseteq M\), let
\[
  \Lang(A)=\Lang\cup\{c_a:a\in A\},
\]
where the \(c_a\) are pairwise distinct new constant symbols, and let
\(\mathcal M_A\) be the \(\Lang(A)\)-expansion of \(\mathcal M\) in which
\(c_a^{\mathcal M_A}=a\).  The notation is needed to distinguish formulas
that may name members of \(A\) from formulas in the original signature.  It
names a signature expansion and an interpreted structure, not a theory, and it
will be used both for diagrams and for types over parameters.
\end{definition}

\begin{definition}[Atomic and elementary diagrams]
\label{def:diagrams}
For a \(\Lang\)-structure \(\mathcal M\), its \emph{atomic diagram}
\(\operatorname{Diag}(\mathcal M)\) is the set of all atomic and negated
atomic \(\Lang(M)\)-sentences true in \(\mathcal M_M\).  Its
\emph{elementary diagram} is
\[
  \operatorname{ElDiag}(\mathcal M)
  =\Th_{\Lang(M)}(\mathcal M_M),
\]
the set of all \(\Lang(M)\)-sentences true in \(\mathcal M_M\).  Both
objects are sets of sentences.  The atomic diagram records exactly the
quantifier-free atomic information needed to test embeddings, whereas the
elementary diagram records every formula needed to test elementary
embeddings; the two names are therefore not interchangeable, and the two
tests below use the distinction.
\end{definition}

\begin{theorem}[Diagram tests for embeddings]
\label{thm:diagram-tests}
Let \(\mathcal M\) and \(\mathcal N\) be \(\Lang\)-structures.
\begin{enumerate}
  \item There is an embedding \(h:\mathcal M\to\mathcal N\) iff
        \(\mathcal N\) has a \(\Lang(M)\)-expansion \(\mathcal N'\) such
        that \(\mathcal N'\entails\operatorname{Diag}(\mathcal M)\).
  \item There is an elementary embedding
        \(h:\mathcal M\to\mathcal N\) iff
        \(\mathcal N\) has a \(\Lang(M)\)-expansion \(\mathcal N'\) such
        that \(\mathcal N'\entails\operatorname{ElDiag}(\mathcal M)\).
\end{enumerate}
\textit{Status: complete proof.}
\end{theorem}

\begin{proof}
Given an embedding \(h\), expand \(\mathcal N\) by interpreting
\(c_a^{\mathcal N'}=h(a)\).  Preservation and reflection of atomic formulas
then make every sentence in \(\operatorname{Diag}(\mathcal M)\) true in
\(\mathcal N'\).  Conversely, suppose such an expansion is given and define
\(h(a)=c_a^{\mathcal N'}\).  For \(a\ne b\), the sentence
\(c_a\ne c_b\) belongs to the atomic diagram, so \(h\) is injective.  The
atomic sentences specifying constants and function values show that \(h\)
preserves those symbols, and the positive and negated relation instances show
that it preserves and reflects relations.  Hence \(h\) is an embedding.

The same expansion constructed from an elementary embedding satisfies every
sentence in the elementary diagram.  Conversely, if
\(\mathcal N'\entails\operatorname{ElDiag}(\mathcal M)\), the map just
defined is an embedding by the first part.  For every formula
\(\varphi(\bar x)\) and tuple \(\bar a\) from \(M\), exactly one of
\(\varphi(\bar c_{\bar a})\) and
\(\neg\varphi(\bar c_{\bar a})\) belongs to the elementary diagram.
Satisfaction of that sentence by \(\mathcal N'\) gives
\[
  \mathcal M\entails\varphi(\bar a)
  \quad\Longleftrightarrow\quad
  \mathcal N\entails\varphi(h\bar a),
\]
so \(h\) is elementary.
\end{proof}

For diagrams and the two embedding tests, see
\cite{Hodges1993,Marker2002}.

\begin{example}[The two diagrams make different tests]
For the inclusion \((\N,<)\subseteq(\Z,<)\), interpret each new constant
\(c_n\) by the integer \(n\).  The resulting expansion of \((\Z,<)\)
satisfies the atomic diagram of \((\N,<)\), in agreement with the fact that
the inclusion is an embedding.  It does not satisfy the elementary diagram:
\(\neg\exists y\,(y<c_0)\) is true in \((\N,<)_{\N}\) and false in the
expanded integer order.
\end{example}

